\section{The original adjacent-pair Hamiltonian}
\label{sec:huang_positive_chain}

The original positive adjacent-pair model of \cite{Landau2015Polynomial}, also used
in \cite[Section~5.4]{Huang2015GroundState}, is a special case of the finite-range
chain model. The following identification preserves its operator exactly.

\begin{lemma}[Distinct adjacent-pair supports]
    \label{thm:huang_adjacent_supports}
    \lean{TNLean.MPS.AreaLaw.nearestNeighborChainSupport_injective,
        TNLean.MPS.AreaLaw.card_nearestNeighborChainSupport}
    \leanok
    The map $i\mapsto\{i,i+1\}$ from the original edge indices
    $0\le i<L-1$ into the nonempty range-one supports is injective. Each such
    support has cardinality two.
\end{lemma}
\begin{proof}
    \leanok
    The least site of the pair is its edge index. Its two consecutive sites are
    distinct and both lie on the chain.
\end{proof}

\begin{definition}[Embedding the positive adjacent-pair interaction]
    \label{def:huang_positive_chain_hamiltonian}
    \lean{TNLean.MPS.AreaLaw.ChainHamiltonian.ofNearestNeighborPositive}
    \leanok
    \uses{thm:huang_adjacent_supports}
    Let $h_i$ be supported on sites $i,i+1$, with $0\le h_i\le I$.
    For each nonempty range-one support $X$, assign the term
    \begin{align}
        \widetilde h_X=
        \begin{cases}
            h_i, & X=\{i,i+1\},      \\
            0,   & \text{otherwise}.
        \end{cases}
    \end{align}
    These terms define a range-one chain Hamiltonian with local norm bound one.
\end{definition}

\begin{lemma}[The interaction terms are preserved]
    \label{thm:huang_positive_chain_terms}
    \lean{TNLean.MPS.AreaLaw.ChainHamiltonian.ofNearestNeighborPositive_term,
        TNLean.MPS.AreaLaw.ChainHamiltonian.ofNearestNeighborPositive_term_eq_zero,
        TNLean.MPS.AreaLaw.ChainHamiltonian.ofNearestNeighborPositive_term_eq_zero_of_card_one,
        TNLean.MPS.AreaLaw.ChainHamiltonian.ofNearestNeighborPositive_term_nonneg,
        TNLean.MPS.AreaLaw.ChainHamiltonian.ofNearestNeighborPositive_term_le_one}
    \leanok
    \uses{def:huang_positive_chain_hamiltonian}
    The term on each designated adjacent pair is the original $h_i$; all
    singleton terms vanish. Every term is supported on its assigned support,
    is Hermitian, and satisfies $0\le\widetilde h_X\le I$ and
    $\|\widetilde h_X\|\le1$.
\end{lemma}
\begin{proof}
    \leanok
    Injectivity makes the assignment well-defined. A zero term has all the
    stated properties. An assigned term inherits its support and positive
    contraction inequalities; these imply Hermiticity and the norm bound.
    A singleton cannot be an adjacent-pair support because their cardinalities
    differ.
\end{proof}

\begin{lemma}[Equality of the Hamiltonian operators]
    \label{thm:huang_positive_chain_operator}
    \lean{TNLean.MPS.AreaLaw.ChainHamiltonian.ofNearestNeighborPositive_operator}
    \leanok
    \uses{thm:huang_positive_chain_terms}
    The constructed finite-range operator is exactly
    $\widetilde H=\sum_X\widetilde h_X=\sum_{i=0}^{L-2}h_i$.
    This includes chains with zero or one site and local dimension zero.
\end{lemma}
\begin{proof}
    \leanok
    Restrict the support sum to the image of the injective edge-support map;
    every term outside this image is zero. Reindexing the remaining finite sum
    by the original edges gives the asserted identity.
\end{proof}
